\section{Encapsulamiento a nivel de LangChain}

\subsection{with\_structured\_output}

El método \texttt{with\_structured\_output} de LangChain simplifica aún más el uso:

\begin{itemize}
    \item Se pasa directamente un Pydantic Model
    \item Internamente se convierte de manera automática a JSON Schema y se coloca en response\_format
    \item El resultado devuelto ya es un objeto instanciado (un paso más de deserialización automática respecto a la API de OpenAI)
\end{itemize}

\begin{warningbox}{Diferencia entre las dos rutas}
\begin{itemize}
    \item \texttt{with\_structured\_output} $\rightarrow$ pasa por \texttt{response\_format} (el JSON Schema se alinea con la semántica de los campos)
    \item \texttt{bind\_tools} $\rightarrow$ pasa por \texttt{tools} (semántica de Function Calling)
\end{itemize}
Un Pydantic Model es, en esencia, también una especie de Function (instanciar la clase $\approx$ llamar al constructor), de modo que bind\_tools de LangChain también puede aceptar un Pydantic Model.
\end{warningbox}

\subsection{Resumen de la sección}

LangChain agrega un encapsulamiento adicional sobre la API de OpenAI: serialización automática, deserialización automática, con una experiencia de uso más simple.

